\documentclass[12pt, a4paper]{scrartcl}
\usepackage[margin=2.5cm]{geometry}
\usepackage[hidelinks]{hyperref}
\usepackage[ngerman]{babel}
\usepackage{graphicx,xcolor,enumitem,booktabs,tabularx,longtable}
\usepackage[T1]{fontenc}
\usepackage[utf8]{inputenc}
\usepackage{lmodern}
\usepackage{parskip}
\usepackage{csquotes}
\usepackage{listings}
\usepackage{caption}
\usepackage{tikz}
\usetikzlibrary{arrows.meta,positioning,shapes.geometric,calc}

% ── Farben ──────────────────────────────────────────────
\definecolor{thblue}{RGB}{0,90,160}
\definecolor{codebg}{RGB}{245,245,245}
\definecolor{codekw}{RGB}{0,90,160}
\definecolor{codecm}{RGB}{100,120,100}
\definecolor{codest}{RGB}{160,40,40}

\lstdefinestyle{code}{
  backgroundcolor=\color{codebg},
  basicstyle=\ttfamily\footnotesize,
  keywordstyle=\color{codekw}\bfseries,
  commentstyle=\color{codecm}\itshape,
  stringstyle=\color{codest},
  breaklines=true,
  showstringspaces=false,
  frame=single,
  framesep=4pt,
  rulecolor=\color{black!20},
  numbers=left,
  numberstyle=\tiny\color{black!50},
  numbersep=8pt,
  tabsize=2,
  captionpos=b,
}
\lstset{style=code}

\captionsetup{labelfont=bf,font=small,skip=4pt}

% ── Metadaten ───────────────────────────────────────────
\newcommand{\faculty}{Fakultät Angewandte Informatik}
\newcommand{\studies}{Bachelor Cyber-Security}
\newcommand{\coursetitle}{Security Engineering}
\newcommand{\semester}{SS\,2026}
\newcommand{\projecttitle}{Secure DevOps (DevSecOps) -- Sicherheitstechnische Betrachtung automatisierter CI/CD-Pipelines}
\newcommand{\supervisor}{Prof.\ Dr.\ Martin Schramm}
\newcommand{\studentA}{Christof Renner (22301943)}
\newcommand{\studentB}{Manuel Friedl (12306626)}
\newcommand{\submissiondate}{\today}

\begin{document}

% ── Titelseite ──────────────────────────────────────────
\begin{titlepage}
  \centering
  \vspace*{1cm}
  {\LARGE Technische Hochschule Deggendorf\\[0.2cm]
   \large \faculty\par}
  \vspace{0.4cm}
  {\Large Studiengang \studies\\[0.3cm]
   Kurs: \coursetitle{} -- \semester\par}
  \vspace{2.5cm}
  {\huge\bfseries\color{thblue} Projektdokumentation\par}
  \vspace{1.5cm}
  {\LARGE\bfseries \projecttitle\par}
  \vfill
  \begin{minipage}[t]{0.48\textwidth}
    \textbf{Vorgelegt von:}\\[6pt]
    \studentA\\
    \studentB
  \end{minipage}\hfill
  \begin{minipage}[t]{0.48\textwidth}\raggedleft
    \textbf{Betreuer:}\\[6pt]
    \supervisor
  \end{minipage}
  \vspace{1.5cm}

  Datum: \submissiondate
\end{titlepage}

\tableofcontents
\newpage

% ────────────────────────────────────────────────────────
\section{Einleitung}

\subsection{Motivation und Problemstellung}
Moderne Softwareentwicklung basiert auf hochgradig automatisierten Build-, Test- und
Deployment-Prozessen. CI/CD-Pipelines sind tief in cloud-native Infrastrukturen
integriert und verfügen über weitreichende Zugriffsrechte auf sensible Ressourcen
wie Quellcode-Repositories, Secrets, Container-Registries und Produktionsumgebungen.
Angriffe auf Softwarelieferketten -- etwa SolarWinds (2020), Codecov (2021),
\mbox{Log4Shell} (2021) und der 3CX-Vorfall (2023) -- belegen, dass der
Entwicklungsprozess selbst zu einem primären Angriffsziel geworden ist.

Aus Sicht des Security Engineering ist eine Pipeline weniger ein Tooling-Problem
als ein \emph{Systemproblem}: Sie verbindet menschliche Identitäten,
externe Abhängigkeiten, Build-Hosts, Registries und Deployment-Targets über
mehrere Vertrauensgrenzen hinweg. Eine isolierte Betrachtung einzelner Werkzeuge
greift daher zu kurz.

\subsection{Zielsetzung}
Ziel der Arbeit ist eine systematische sicherheitstechnische Betrachtung einer
repräsentativen CI/CD-Pipeline. Konkret werden:
\begin{itemize}[leftmargin=1.4em]
  \item schützenswerte Assets und Trust Boundaries identifiziert,
  \item Bedrohungen je Pipeline-Stufe nach STRIDE strukturiert analysiert,
  \item geeignete Sicherheitsmechanismen abgeleitet und in einer
        prototypischen Pipeline (GitHub Actions) umgesetzt,
  \item die Wirksamkeit dieser Maßnahmen anhand von SLSA-Levels und
        OpenSSF-Scorecard evaluiert,
  \item verbleibende Restrisiken transparent gemacht.
\end{itemize}

\subsection{Abgrenzung}
Nicht Gegenstand der Arbeit sind: vollständige Härtung des Cloud-Tenants,
Penetrationstests der Pipeline-Plattform selbst sowie Implementierung organisatorischer
Maßnahmen jenseits der Branch-Protection-Konfiguration.

\subsection{Aufbau der Arbeit}
Kapitel~\ref{sec:grundlagen} klärt die Begriffe DevSecOps, Supply-Chain-Security
und die zugrundeliegenden Frameworks (STRIDE, SLSA, SSDF). Kapitel~\ref{sec:system}
beschreibt das modellierte System und seine Trust Boundaries.
Kapitel~\ref{sec:threatmodel} enthält das STRIDE-Threat-Model.
Kapitel~\ref{sec:countermeasures} ordnet jeder Bedrohung konkrete
Gegenmaßnahmen zu und dokumentiert die Implementierung im Prototyp.
Kapitel~\ref{sec:evaluation} bewertet die Wirksamkeit, Kapitel~\ref{sec:fazit}
fasst zusammen und diskutiert Restrisiken.

% ────────────────────────────────────────────────────────
\section{Grundlagen}\label{sec:grundlagen}

\subsection{DevSecOps}
DevSecOps bezeichnet die Integration sicherheitstechnischer Aktivitäten in
den DevOps-Lebenszyklus mit dem Ziel, Sicherheit als kontinuierlichen,
automatisierten Bestandteil von Build und Deployment zu etablieren. 
Sicherheitsmaßnahmen werden nicht nachgelagert aufgesetzt, sondern als Teil derselben Pipeline
ausgeführt, die auch das Artefakt produziert.

\subsection{Software Supply Chain Security}
Eine Software Supply Chain umfasst alle Personen, Prozesse, Werkzeuge und
Abhängigkeiten, die zur Erstellung und Auslieferung eines Artefakts beitragen.
Angriffsklassen nach \cite{slsa} sind insbesondere:
\begin{itemize}[leftmargin=1.4em]
  \item \textbf{Source}-Angriffe (z.\,B.\ kompromittierte Maintainer-Accounts,
        Typosquatting),
  \item \textbf{Build}-Angriffe (z.\,B.\ kompromittierte Runner, manipulierte
        Compiler/Plugins),
  \item \textbf{Dependency}-Angriffe (z.\,B.\ Übernahme verwaister Pakete),
  \item \textbf{Distribution}-Angriffe (z.\,B.\ Image-Tag-Replacement in
        Registries).
\end{itemize}

\subsection{Frameworks}
\begin{description}[leftmargin=1.4em,style=nextline]
  \item[STRIDE] Microsoft-Methodik zur Bedrohungsmodellierung mit den
        Kategorien \emph{Spoofing, Tampering, Repudiation, Information
        Disclosure, Denial of Service, Elevation of Privilege}.
  \item[SLSA v1.0] \emph{Supply-chain Levels for Software Artifacts}
        definiert Anforderungen an Build-Integrität in vier Stufen
        (Build~L1--L3, Source~L1--L3).
  \item[NIST SSDF (SP\,800-218)] \emph{Secure Software Development Framework}
        -- Praktiken über den gesamten SDLC.
  \item[OpenSSF Scorecard] Automatisierte Bewertung von Repository-Posture
        anhand objektiver Checks.
\end{description}

\subsection{Security Engineering Framework}\label{sec:seframework}
Die Arbeit folgt dem in der Lehrveranstaltung eingeführten
Vier-Quadranten-Modell des Security Engineering nach Anderson:
\emph{Policy}, \emph{Mechanism}, \emph{Assurance} und \emph{Incentives}.
Tabelle~\ref{tab:seframework} ordnet jedem Quadranten konkrete
Bestandteile der untersuchten Pipeline zu; auf den Quadranten
\emph{Incentives} wird in Abschnitt~\ref{sec:humanfactor} vertiefend
eingegangen.

\begin{table}[h]
\centering
\begin{tabular}{p{2.4cm} p{11.5cm}}
\toprule
\textbf{Quadrant} & \textbf{Umsetzung in der Pipeline} \\
\midrule
Policy &
Branch-Protection-Regeln, CODEOWNERS, Required Status Checks,
Coverage-Schwelle 80\,\%, HIGH/CRITICAL als Build-Breaker,
Manual-Approval-Gate für Production. \\
Mechanism &
Bandit, CodeQL, pip-audit, Trivy, Syft (SBOM), Cosign Keyless Signing,
SLSA-Provenance, OIDC-Federation, Action-Pinning auf Commit-SHA. \\
Assurance &
SLSA-Level-Bewertung, OpenSSF Scorecard (wöchentlich),
SARIF-Upload in Code-Scanning, Cosign-Verify im CD-Workflow,
Sigstore Rekor Transparency Log. \\
Incentives &
Vier-Augen-Prinzip via CODEOWNERS, automatisierte Dependabot-PRs
reduzieren Reibung für sichere Updates, parallele Jobs +
\texttt{cancel-in-progress} mildern den Geschwindigkeits-Konflikt
(siehe Abschnitt~\ref{sec:humanfactor}). \\
\bottomrule
\end{tabular}
\caption{Mapping der Pipeline-Bestandteile auf das
Security-Engineering-Framework.}
\label{tab:seframework}
\end{table}

% ────────────────────────────────────────────────────────
\section{System unter Betrachtung}\label{sec:system}

\subsection{Architekturüberblick}
Untersucht wird eine prototypische, container-basierte Pipeline für eine
Python-Webanwendung. Die wesentlichen Stufen sind in
Abbildung~\ref{fig:pipeline} dargestellt.

\begin{figure}[h]
\centering
\begin{tikzpicture}[
  node distance=0.5cm and 0.7cm,
  stage/.style={rectangle, rounded corners, draw=thblue, thick, fill=thblue!8,
                minimum width=2.4cm, minimum height=1.0cm, align=center,
                font=\small},
  arrow/.style={-Stealth, thick, thblue},
  ext/.style={rectangle, draw=black!50, dashed, fill=black!5,
              minimum width=2.4cm, minimum height=0.9cm, align=center,
              font=\footnotesize\itshape}
]
  \node[stage] (src)    {Source\\(GitHub)};
  \node[stage, right=of src]   (build)  {Build\\(Runner)};
  \node[stage, right=of build] (test)   {Test \&\\Scan};
  \node[stage, right=of test]  (pkg)    {Package\\\& Sign};
  \node[stage, right=of pkg]   (deploy) {Deploy\\(Cloud)};

  \draw[arrow] (src)   -- (build);
  \draw[arrow] (build) -- (test);
  \draw[arrow] (test)  -- (pkg);
  \draw[arrow] (pkg)   -- (deploy);

  \node[ext, below=1.0cm of build] (deps) {PyPI / Base\\Images};
  \node[ext, below=1.0cm of pkg]   (reg)  {GHCR \&\\Sigstore};

  \draw[arrow, dashed] (deps) -- (build);
  \draw[arrow, dashed] (pkg)  -- (reg);
\end{tikzpicture}
\caption{Pipeline-Stufen und externe Vertrauensanker.}
\label{fig:pipeline}
\end{figure}

\subsection{Schützenswerte Assets}
\begin{longtable}{p{4.2cm} p{9.5cm}}
\toprule
\textbf{Asset} & \textbf{Schutzziel (CIA + AAA)} \\
\midrule \endhead
Quellcode (main) & Integrität, Authentizität \\
Build-Artefakt (Image) & Integrität, Authentizität, Verfügbarkeit \\
Secrets (Credentials, Tokens) & Vertraulichkeit, Integrität \\
SBOM \& Provenance & Integrität, Nicht-Abstreitbarkeit \\
Pipeline-Konfiguration & Integrität, Autorisierung \\
Build-Runner-Identität (OIDC) & Authentizität, Vertraulichkeit \\
\bottomrule
\end{longtable}

\subsection{Trust Boundaries}
\begin{enumerate}[leftmargin=1.4em]
  \item Entwickler-Workstation $\leftrightarrow$ GitHub
  \item GitHub-Plattform $\leftrightarrow$ Build-Runner
  \item Build-Runner $\leftrightarrow$ Paket-/Image-Registries (PyPI, Docker Hub, GHCR)
  \item Build-Runner $\leftrightarrow$ Sigstore (Fulcio/Rekor)
  \item Pipeline $\leftrightarrow$ Deployment-Target (Cloud-API)
\end{enumerate}

\subsection{Angreifermodell}\label{sec:adversary}
Die Lehrveranstaltung kategorisiert relevante Gegnertypen entlang ihrer
Motivation, Ressourcen und Methoden: Cyberkriminelle, staatliche
Akteure (Nachrichtendienste), unehrliche Insider, konkurrierende
Unternehmen sowie -- im Sinne Andersons -- aktivistisch motivierte
Gruppen (\enquote{wütender Mob}). Diese Kategorien werden für die
vorliegende Arbeit auf drei DevSecOps-relevante Profile verdichtet,
die entlang der Trust Boundaries (Abschnitt 3.3) operieren:

\begin{description}[leftmargin=1.4em,style=nextline]
  \item[A1 -- Externer Angreifer] kein Repo-Zugriff; nutzt Phishing,
        Typosquatting, Schwachstellen in Abhängigkeiten und
        Tag-Hijacks in öffentlichen Registries.
        Deckt die Vorlesungs-Kategorien
        \emph{Cyberkriminelle} (finanziell motiviert: Crypto-Mining auf
        Runnern, Ransomware) und \emph{Hacktivismus} (z.\,B.\
        Protestware in npm-Paketen) ab.
  \item[A2 -- Kompromittierter Mitwirkender] gültige, aber begrenzte
        Schreibrechte -- Fork-PR-Author, einzelner Maintainer mit
        übernommenem Account, oder ein Maintainer einer als Dependency
        eingezogenen Bibliothek. Entspricht der Vorlesungskategorie
        \emph{unehrliche Insider} und ist zugleich der typische
        Eintrittsvektor für staatliche Supply-Chain-Operationen
        (SolarWinds-Stil, vgl.\ Walkthrough~1 in
        Abschnitt~\ref{sec:walkthroughs}).
  \item[A3 -- Insider/Plattform] Admin-Rechte am Repo, kompromittierter
        Self-hosted-Runner, oder ein staatlicher Akteur mit
        Zugriff auf einen Vertrauensanker (Registry, Sigstore, GitHub).
        Vereint die Vorlesungs-Kategorien \emph{Nachrichtendienste}
        (ressourcenstark, verdeckt) und \emph{Insider mit
        Admin-Rechten}; in beiden Fällen ist der Angreifer
        innerhalb des deklarierten Vertrauensbereichs der Pipeline.
\end{description}

Die Differenzierung in drei Profile ist bewusst schmal gehalten:
Eine feinere Aufschlüsselung würde die folgende STRIDE-Analyse
(Tab.~\ref{tab:stride}) künstlich aufblähen, ohne neue
Mitigationsentscheidungen sichtbar zu machen. Die Vorlesungs-Kategorie
\emph{konkurrierende Unternehmen} wird hier nicht eigenständig
modelliert, da deren Angriffsvektoren technisch deckungsgleich mit
A1 oder A2 sind und sich lediglich in der Motivation unterscheiden.

% ────────────────────────────────────────────────────────
\section{Bedrohungsanalyse (STRIDE)}\label{sec:threatmodel}

\subsection{STRIDE-Mapping pro Pipeline-Stufe}
Die folgende Tabelle dokumentiert pro Pipeline-Stufe die wesentlichen
Bedrohungen, klassifiziert nach STRIDE. Die Risiko-Einstufung folgt einer
qualitativen Skala (niedrig/mittel/hoch) auf Basis von Eintrittswahrscheinlichkeit
und Schadensausmaß im Kontext einer Studienarbeit-typischen Demo-Pipeline.

\begin{longtable}{p{1.6cm} p{2.0cm} p{7.4cm} p{1.4cm} p{1.4cm}}
\toprule
\textbf{Stufe} & \textbf{STRIDE} & \textbf{Bedrohung} & \textbf{Angr.} & \textbf{Risiko} \\
\midrule \endhead
Source & S & Push mit gestohlenen Maintainer-Credentials & A1/A2 & hoch \\
Source & T & Manipulation von Pipeline-Definitionen via PR & A2 & hoch \\
Source & R & Unsignierte Commits, fehlende Audit-Trails & A2 & mittel \\
Source & I & Secrets im Klartext im Repo (Logs, Configs) & A1 & hoch \\
Build  & T & Bösartige transitive Dependency (Confused-Deps, Typosquat) & A1 & hoch \\
Build  & T & Manipulation des Base-Images (Tag-Replacement) & A1 & hoch \\
Build  & E & Workflow mit zu weiten \texttt{permissions} (RW everywhere) & A2 & hoch \\
Build  & E & \texttt{pull\_request\_target} mit \texttt{secrets} ausgesetzt & A1 & hoch \\
Build  & I & Cache-Poisoning zwischen PRs & A2 & mittel \\
Test   & D & Lange laufende Scans blockieren Deployments & A1 & niedrig \\
Test   & T & Test-Bypass via geänderte Workflow-Bedingungen & A2 & mittel \\
Pkg    & T & Ungesignedes Image, Tag-Hijack in Registry & A1 & hoch \\
Pkg    & R & Fehlende Provenance / SBOM & -- & mittel \\
Deploy & S & Long-lived Cloud-Credentials in Repo-Secrets & A2 & hoch \\
Deploy & E & Fehlende Approval-Gates für Production & A2 & hoch \\
Deploy & I & Verbose Logs leaken Tokens / PII & A1 & mittel \\
\bottomrule
\caption{STRIDE-Mapping der untersuchten Pipeline.}
\label{tab:stride}
\end{longtable}

\newpage

\subsection{Threat Walkthroughs}\label{sec:walkthroughs}
Die kondensierte Tabellenform aus Tabelle~\ref{tab:stride} verdichtet
Bedrohungen, verschleiert aber ihre tatsächliche \emph{Verkettung}. Reale
Angriffe auf CI/CD-Systeme bestehen selten aus einem einzelnen Schritt;
charakteristisch ist die schrittweise Eskalation entlang der Trust
Boundaries. Im Folgenden werden drei repräsentative Bedrohungen aus
Tabelle~\ref{tab:stride} als ausformulierte Angriffsketten dargestellt --
inklusive Voraussetzungen, Eskalationsschritten, Erkennungs- und
Mitigationspunkten.

\paragraph{Walkthrough 1 -- Bösartige transitive Dependency
(Build-T, A1).}
Eine populäre, indirekt eingezogene Abhängigkeit wird durch einen
kompromittierten Maintainer-Account aktualisiert; die neue Version
enthält eine Post-Install-Routine, die Umgebungsvariablen und
\texttt{\textasciitilde/.aws/credentials} an einen externen
Endpunkt sendet. Dependabot erstellt plangemäß einen Update-PR;
Reviewer prüfen den Diff der \texttt{requirements.txt} und sehen
nur einen scheinbar harmlosen Versionssprung der direkten
Abhängigkeit -- die transitive Aktualisierung ist nicht sichtbar.
Beim CI-Build läuft die Post-Install-Routine mit den
Workflow-Permissions und exfiltriert \texttt{\$GITHUB\_TOKEN} sowie
gemountete Cloud-Credentials.

\textit{Erkennung und Mitigation.} pip-audit greift bei einer
frischen Kompromittierung nicht (nur bekannte CVEs); zweite
Verteidigungslinie ist der SBOM-Diff
(Abschnitt~\ref{sec:sbomlifecycle}), Egress-Restriktionen können
die Exfiltration zusätzlich blockieren.
\texttt{permissions: \{\}} reduziert die Token-Reichweite,
\texttt{persist-credentials: false} verhindert Token-Persistenz im
\texttt{.git/config}, OIDC-Federation entzieht dem Angreifer
langlebige Cloud-Secrets. Restrisiko: Der Angriff bleibt möglich,
solange transitive Updates nicht an eine kuratierte Allowlist
gebunden sind -- weitergehend wäre ein Privat-Mirror mit manueller
Freigabe.

\paragraph{Walkthrough 2 -- Manipulation der Pipeline-Definition via PR
(Source-T, A2).}
Ein Kontributor mit \enquote{Triage}-Rolle reicht einen
unscheinbaren Pull-Request ein, der nicht den Anwendungscode,
sondern eine Datei unter \texttt{.github/workflows/} ändert -- etwa
um einen zusätzlichen Schritt einzufügen, der
\texttt{secrets.AWS\_DEPLOY\_KEY} in einen externen Webhook
schreibt. Optisch ist der PR ein \enquote{Refactoring}: Der
Workflow wird in mehrere Teildateien zerlegt, dabei werden
zusätzliche \texttt{run:}-Schritte eingefügt. Reviewer ohne
explizit gepflegte CODEOWNERS-Regeln neigen dazu, den Diff zu
überfliegen; sobald der PR gemergt ist und der Workflow auf
\texttt{main} läuft, greift der eingefügte Schritt auf die vollen
\texttt{secrets.\textasciitilde}-Werte zu.

\textit{Erkennung und Mitigation.} CODEOWNERS auf \texttt{/.github/}
verschiebt den Review zwingend an die Pipeline-Verantwortlichen;
OpenSSF Scorecard markiert das Fehlen dieser Regel als
Posture-Defizit, \texttt{actionlint} und \texttt{zizmor} prüfen in
der CI auf typische PPE-Antipatterns. Required-Status-Checks für
jede Pipeline-Änderung sowie die Regel, dass \texttt{secrets} nur in
Jobs eines Approval-Environments verfügbar sind, schließen den
Vektor weitgehend. Restrisiko: Ein Insider mit Admin-Rechten kann
CODEOWNERS selbst ändern -- hier hilft nur Org-weites Audit-Logging
mit Alarm bei Branch-Protection-Änderungen.

\paragraph{Walkthrough 3 -- Tag-Hijack in der Registry
(Pkg-T, A1).}
Der Angreifer hat sich Schreibzugriff auf den Registry-Namespace
verschafft -- z.\,B.\ über einen kompromittierten
Service-Account-Token, der versehentlich in einem öffentlichen
Issue gepostet wurde. Er baut lokal ein bösartiges Image, vergibt
denselben Tag wie das letzte Release
(\texttt{ghcr.io/org/app:1.4.2}) und pusht es; der Tag zeigt jetzt
auf das neue Manifest. Ein Deployment-Skript, das nur den Tag
referenziert, zieht beim nächsten Rollout das manipulierte Image.

\textit{Erkennung und Mitigation.} Der CD-Workflow verifiziert vor
dem Rollout per \texttt{cosign verify} sowohl die Signatur als auch
die SLSA-Provenance-Attestation gegen die OIDC-Identität des
CI-Workflows. Beide Checks schlagen fehl, da der Angreifer keine
gültige Signatur unter dem CI-Workflow-OIDC-Issuer erzeugen kann;
das Rekor-Transparency-Log macht Manipulationen zusätzlich
öffentlich auditierbar. Deploy-Schritte referenzieren ausschließlich
den Image-Digest (\texttt{@sha256:\dots}), nicht den Tag. Restrisiko:
Ein Angriff auf Sigstore selbst ist nicht unmittelbar mitigierbar --
Sigstore ist daher in Abschnitt~\ref{sec:system} explizit als
Vertrauensanker geführt.

\subsection{Vertiefung: Poisoned Pipeline Execution (PPE)}
\label{sec:ppe}
Die Bedrohungen \emph{Build-T} (Manipulation Pipeline-Definition) und
\emph{Build-E} (\texttt{pull\_request\_target} mit Secrets) aus
Tabelle~\ref{tab:stride} sind manifestationen einer übergeordneten
Klasse, die OWASP als \emph{Poisoned Pipeline Execution}
\cite{owaspcicd} beschreibt. PPE bezeichnet jeden Angriff, der den
\enquote{Konfigurations-als-Code}-Charakter moderner CI-Systeme
ausnutzt: Wer den Workflow ändern darf, kontrolliert -- transitiv --
alle Aktionen, die in seinem Kontext ablaufen.

\paragraph{Direkte und indirekte PPE.}
Direkte PPE liegt vor, wenn der Angreifer die Pipeline-Definition
selbst modifiziert (z.\,B.\ als Maintainer mit Push-Recht). Indirekte
PPE nutzt Trigger aus, die \emph{trotz} fehlender Schreibrechte
Code im Kontext der Pipeline ausführen.

\paragraph{Anatomie von \texttt{pull\_request\_target}.}
Im Gegensatz zum Standard-Trigger \texttt{pull\_request} -- der den
Workflow im Kontext des PR-Branches ohne Zugriff auf
\texttt{secrets} ausführt -- läuft \texttt{pull\_request\_target}
im Kontext des Ziel-Branches \emph{mit} vollem Secret-Zugriff. Der
Trigger wurde geschaffen, um Aktionen wie automatisches Labeling
auch für Forks zu ermöglichen, bei denen der Forker keinen
Schreibzugriff auf das Ziel-Repository besitzt. Genau diese Kombination
(Code-Quelle = Fork, Secrets = Original-Repo) ist die Kernschwäche.

\begin{lstlisting}[language=bash,caption={Anti-Pattern: \texttt{pull\_request\_target} mit unsicherem Checkout},label={lst:ppe}]
on:
  pull_request_target:        # laeuft im Ziel-Repo-Kontext mit Secrets
    types: [opened, synchronize]
jobs:
  ci:
    steps:
      - uses: actions/checkout@<sha>   # OK: HEAD = Ziel-Branch
        with:
          ref: ${{ github.event.pull_request.head.sha }}  # PROBLEM!
      - run: npm ci && npm test  # PROBLEM: laedt + fuehrt Fork-Code aus
\end{lstlisting}

Ein bösartiger Forker fügt seinem PR-Branch einen
\texttt{postinstall}-Hook in der \texttt{package.json} hinzu, der die
Werte aller Umgebungsvariablen exfiltriert. Beim ersten
\texttt{npm ci} im Workflow läuft dieser Hook -- inklusive Zugriff auf
\texttt{secrets.NPM\_TOKEN}, \texttt{secrets.AWS\_*} und das
\texttt{GITHUB\_TOKEN} mit Schreibrecht auf das Original-Repo.

\paragraph{Codecov 2021 als Realbeispiel.}
Im April 2021 wurde der Codecov-Bash-Uploader durch einen Fehler in
einem Docker-Image-Build kompromittiert; in der Folge konnten
Angreifer rund zwei Monate lang Umgebungsvariablen aus
CI-Umgebungen exfiltrieren, die den Uploader integriert hatten. Der
Vorfall war zwar kein klassischer PPE-Fall, illustriert aber die
gleiche Kausalkette: ein in einem Pipeline-Schritt ausgeführtes
Skript, das implizit Zugriff auf alle Secrets erhält, die im Kontext
des Workflows verfügbar sind. Branchenweite Reaktion war eine
Auseinandersetzung mit dem Trigger-Modell und der
Secret-Sichtbarkeit.

\paragraph{Mitigation in der hier umgesetzten Pipeline.}
Die folgenden konstruktiven Entscheidungen schließen die
PPE-Klasse weitgehend:
\begin{itemize}[leftmargin=1.4em]
  \item \textbf{Kein \texttt{pull\_request\_target}.} Der CI-Workflow
        verwendet ausschließlich den sicheren \texttt{pull\_request}-Trigger;
        Aktionen mit Secret-Bedarf (z.\,B.\ Image-Push) werden bei PRs
        per \texttt{if: github.event\_name != 'pull\_request'}
        ausgesperrt.
  \item \textbf{CODEOWNERS auf \texttt{/.github/}.} Jede
        Workflow-Änderung erfordert ein Review der designierten
        Pipeline-Verantwortlichen; das Vier-Augen-Prinzip wird damit auf
        die kritischste Datei-Klasse angewendet.
  \item \textbf{Default-deny \texttt{permissions: \{\}}.} Selbst
        wenn ein PPE-Angriff Code im Workflow-Kontext zur Ausführung
        bringt, ist das verfügbare \texttt{GITHUB\_TOKEN} ohne
        Schreibrechte und kann weder Branches noch Releases
        manipulieren.
  \item \textbf{OIDC-Federation statt langlebiger Secrets.}
        Cloud-Zugriffe verwenden kurzlebige, an die OIDC-Identität des
        Workflows gebundene Tokens. Selbst eine erfolgreiche
        Exfiltration nutzt dem Angreifer wenig, da die Tokens auf
        Minutenbasis ablaufen und an die Workflow-Run-ID gebunden sind.
  \item \textbf{\texttt{persist-credentials: false}.} Der
        \texttt{GITHUB\_TOKEN} wird nach \texttt{actions/checkout}
        nicht in \texttt{.git/config} verewigt; ein nachgelagertes
        Skript hat damit keinen Hebel mehr.
\end{itemize}

\paragraph{Verbleibende PPE-Risiken.}
Eine vollständige Immunität ist nicht erreichbar. Eine
Kompromittierung der GitHub-Plattform selbst, eine Schwachstelle in
einer als Dependency genutzten Action (vgl.\ Fallstudie
\texttt{tj-actions} in Abschnitt~\ref{sec:tjactions}) oder ein
Insider mit Administrationsrechten am Repo bleiben außerhalb der
Reichweite der hier diskutierten Maßnahmen. Diese Risiken werden in
Abschnitt~\ref{sec:restrisks} explizit als Restrisiken geführt.

% ────────────────────────────────────────────────────────
\section{Gegenmaßnahmen und Implementierung}\label{sec:countermeasures}

\subsection{Source-Stufe}
\begin{description}[leftmargin=1.4em,style=nextline]
  \item[Branch Protection \& CODEOWNERS] \texttt{main} ist geschützt; Merges
        erfordern ein Review durch Code-Owner (Vier-Augen-Prinzip), bestandene
        Pflicht-Checks, signierte Commits und lineare History.
  \item[Signierte Commits] Entwickler signieren mit SSH/GPG; nicht-signierte
        Commits werden serverseitig abgelehnt.
  \item[Secret-Scanning] \texttt{gitleaks} läuft als Pre-Commit-Hook und in der
        CI; GitHub Push-Protection ergänzt das.
  \item[Pinned Actions] Alle GitHub-Actions sind auf Commit-SHA gepinnt, um
        Tag-Retag-Angriffe (z.\,B.\ \texttt{tj-actions/changed-files} 03/2025)
        zu verhindern.
\end{description}

\subsection{Build-Stufe}
\begin{description}[leftmargin=1.4em,style=nextline]
  \item[Least-Privilege-Permissions] Workflow-Default ist \texttt{permissions: \{\}};
        jeder Job opt-in mit minimalen Scopes.
  \item[Ephemere Runner] Nutzung von GitHub-hosted Runnern (frische VM pro Job);
        Self-hosted bewusst nicht verwendet.
  \item[Reproduzierbares Image] Multi-stage \texttt{Dockerfile} mit gepinntem
        Base-Image, Non-Root-User (UID~10001), Read-only-Filesystem-tauglich.
  \item[Dependency-Pinning] \texttt{requirements.txt} mit exakten Versionen;
        Updates über Dependabot-PRs, die die volle Pipeline durchlaufen.
\end{description}

\subsection{Test- und Scan-Stufe}
\begin{itemize}[leftmargin=1.4em]
  \item \textbf{SAST}: Bandit (Python) + GitHub CodeQL, Ergebnisse als SARIF in
        die GitHub Code-Scanning-API.
  \item \textbf{SCA}: \texttt{pip-audit} gegen die OSV-Datenbank, \texttt{--strict}
        bricht den Build bei bekannter CVE ab.
  \item \textbf{Container-Scan}: Trivy auf das gebaute Image,
        \texttt{HIGH,CRITICAL} mit \texttt{exit-code: 1}.
  \item \textbf{Coverage-Gate}: \texttt{pytest --cov-fail-under=80}.
\end{itemize}

\subsection{Package- und Sign-Stufe}
\begin{description}[leftmargin=1.4em,style=nextline]
  \item[SBOM] Syft erzeugt CycloneDX-JSON aus dem Image; das SBOM wird als
        Attestation an das Image angeheftet.
  \item[Cosign Keyless Signing] Signatur via Sigstore Fulcio mit dem
        OIDC-Token des Workflows; keine langlebigen Schlüssel im Repo.
  \item[SLSA-Provenance] \texttt{actions/attest-build-provenance} erzeugt eine
        verifizierbare Provenance-Attestation (SLSA v1).
\end{description}

\subsection{SBOM-Lebenszyklus}\label{sec:sbomlifecycle}
Eine \emph{Software Bill of Materials} (SBOM) listet alle direkten und
transitiven Komponenten eines Artefakts auf. Sie ist kein einzelner
Sicherheitsmechanismus, sondern der zentrale \emph{Datensatz}, auf dem
mehrere Mechanismen aufsetzen: Schwachstellen-Diff, Lizenz-Tracking,
Compliance-Nachweise und langfristige Reproduzierbarkeit. Die
Implementierung folgt dem in Abbildung~\ref{fig:sbom} skizzierten
Lebenszyklus.

\begin{figure}[h]
\centering
\begin{tikzpicture}[
  node distance=0.45cm and 0.55cm,
  sbomstep/.style={rectangle, rounded corners, draw=thblue, thick,
               fill=thblue!8, minimum width=2.6cm, minimum height=0.9cm,
               align=center, font=\footnotesize},
  arrow/.style={-Stealth, thick, thblue}
]
  \node[sbomstep] (gen)    {Generate\\(Syft)};
  \node[sbomstep, right=of gen]   (att)   {Attest\\(Cosign)};
  \node[sbomstep, right=of att]   (push)  {Push\\(Registry)};
  \node[sbomstep, right=of push]  (verify){Verify\\(CD)};
  \node[sbomstep, right=of verify](diff)  {Diff /\\Monitor};

  \draw[arrow] (gen)    -- (att);
  \draw[arrow] (att)    -- (push);
  \draw[arrow] (push)   -- (verify);
  \draw[arrow] (verify) -- (diff);
\end{tikzpicture}
\caption{SBOM-Lebenszyklus: Erzeugung, Attestation, Verteilung,
Verifikation, kontinuierliches Monitoring.}
\label{fig:sbom}
\end{figure}

\paragraph{1.~Erzeugung mit Syft.}
Syft (Anchore) inspiziert das gebaute Image schichtweise und erkennt
Pakete von Distribution-Manager (\texttt{apt}), Sprach-Manager
(\texttt{pip}, \texttt{npm}, \texttt{gem}) sowie statisch eingebundene
Bibliotheken. Ausgabeformat ist CycloneDX in JSON, weil dies zusätzlich
zur reinen Komponentenliste auch \emph{Beziehungen} (\texttt{depends-on},
\texttt{contains}) und Hashes erfasst -- beides ist für spätere
Verifikation entscheidend.

\paragraph{2.~Attestation mit Cosign.}
Statt das SBOM lose neben dem Image zu lagern (klassisches
\texttt{sbom.json} im Release-Asset-Ordner), wird es als signierte
\emph{Attestation} an den Image-Digest gebunden:
\texttt{cosign attest --predicate sbom.cdx.json --type cyclonedx}.
Damit ist das SBOM kryptographisch an genau das Image gekettet, das
es beschreibt -- und manipulierte SBOM-Kopien lassen sich erkennen.

\paragraph{3.~Verteilung über die Registry.}
Cosign legt Attestations als zusätzliche OCI-Manifeste neben dem
eigentlichen Image ab; Pull-Clients holen sie auf Wunsch
(\texttt{cosign download attestation}). Vorteile: keine zweite
Distributionspipeline, keine zusätzlichen Zugriffsrechte, keine
abweichende Aufbewahrungsdauer.

\paragraph{4.~Verifikation vor Deploy.}
Der CD-Workflow ruft \texttt{cosign verify-attestation} mit
\texttt{--type cyclonedx} und \texttt{--certificate-identity-regexp}
auf, der nur Attestations der erwarteten CI-Workflow-Identität
akzeptiert. Damit wird eine angeheftete SBOM-Manipulation durch einen
externen Angreifer mit Registry-Zugriff zur Pipeline-Stufe
\emph{vor} dem Rollout sichtbar.

\paragraph{5.~Kontinuierliches Monitoring (SBOM-Diff).}
Der eigentliche operative Wert eines SBOM entsteht erst nach
Auslieferung. Bei jedem Build vergleicht ein nachgelagerter Schritt
die neue SBOM gegen die letzte gepushte Version desselben
Stream-Tags. Drei Fragen werden beantwortet:
\begin{enumerate}[leftmargin=1.4em]
  \item \textbf{Welche Pakete sind hinzugekommen?} Neue transitive
        Abhängigkeiten sind das primäre Eintrittsfenster für
        Supply-Chain-Angriffe (vgl.\ Walkthrough~1 in
        Abschnitt~\ref{sec:walkthroughs}).
  \item \textbf{Welche Versionen haben sich geändert?} Ein
        Versionssprung über mehrere Major-Releases hinweg ist ein
        Signal für übersehene Major-Updates und sollte einen Reviewer
        zwingen, das Changelog zu prüfen.
  \item \textbf{Sind bekannte CVEs hinzugekommen?} Werkzeuge wie
        \texttt{grype} oder \texttt{trivy sbom} korrelieren die SBOM
        gegen die OSV/NVD-Datenbanken; das ist günstiger als ein
        kompletter Image-Scan und lässt sich gegen \emph{historische}
        Artefakte ausführen, die in Produktion noch laufen.
\end{enumerate}

\paragraph{6.~Aufbewahrung und Auskunft.}
SBOMs müssen über die Lebensdauer des Artefakts verfügbar bleiben --
für eine Library, die sieben Jahre verkauft wird, also potenziell
deutlich länger als ein typischer Container-Lebenszyklus
(vgl.\ die in Abschnitt~\ref{sec:grundlagen} erwähnte
Software-Sustainability-Diskussion). Die Bindung der Attestation an
den Image-Digest und die Hinterlegung im Sigstore-Rekor-Log adressiert
das auch dann, wenn die ursprüngliche Pipeline-Konfiguration nicht
mehr existiert.

\paragraph{Bezug zum Cyber Resilience Act (CRA).}
Mit dem \emph{Cyber Resilience Act}~\cite{cra} der EU werden SBOMs
ab Ende 2027 für \enquote{Produkte mit digitalen Elementen} eine
verpflichtende technische Dokumentation. Hersteller müssen eine SBOM
für jedes vermarktete Produkt führen und Behörden auf Anfrage
verfügbar machen; bei \enquote{important} bzw.\ \enquote{critical}
Produkten kommen erweiterte Pflichten hinzu (Conformity Assessment,
Vulnerability Handling, Coordinated Disclosure). Die hier umgesetzte
Pipeline erfüllt die kommenden technischen Anforderungen \emph{by
construction}: SBOM wird automatisiert erzeugt, an das Artefakt
gebunden, signiert verteilt und ist über das Rekor-Log
nachvollziehbar. Die organisatorischen CRA-Pflichten (z.\,B.\
24h-Reporting bei aktiv ausgenutzten Schwachstellen) bleiben außerhalb
des Pipeline-Scopes, sind aber durch die SECURITY.md-Policy und das
Vulnerability-Disclosure-Verfahren des Repositories vorbereitet.

\paragraph{Grenzen.}
Eine SBOM ist nur so vollständig wie der Generator. Statisch
eingebettete C-Bibliotheken in Wheel-Distributionen,
vendoring-basierte Build-Systeme und nachgeladene Plugins entgehen
einer rein paketmanager-basierten Erkennung.

\subsection{Deploy-Stufe}
\begin{description}[leftmargin=1.4em,style=nextline]
  \item[OIDC-Federation] Cloud-Logins (Azure / AWS) erfolgen via OIDC --
        keine statischen Cloud-Credentials in Repo-Secrets.
  \item[Environments mit Approval] \texttt{production} ist GitHub-Environment
        mit Pflicht-Reviewer und Wait-Timer.
  \item[Cosign-Verify vor Deploy] Vor dem Rollout verifiziert der CD-Workflow
        Signatur \emph{und} SLSA-Provenance des Image-Digests.
  \item[Image-Pinning per Digest] Deployments referenzieren \texttt{@sha256:\dots},
        nicht Tags.
\end{description}

\subsection{Auszug aus der CI-Implementierung}
Listing~\ref{lst:ci} zeigt die Job-Topologie und die zentralen
Sicherheits-Gates.

\begin{lstlisting}[language=bash,caption={Auszug \texttt{.github/workflows/ci.yml}: Permissions \& Pin auf SHA},label={lst:ci}]
permissions: {}              # default-deny

jobs:
  lint-test:
    permissions: { contents: read }
    steps:
      - uses: actions/checkout@11bd71901bbe5b1630ceea73d27597364c9af683 # v4.2.2
        with: { persist-credentials: false }
  build-image:
    needs: [lint-test, sast, sca, secret-scan]
    permissions:
      contents: read
      packages: write
      id-token: write       # OIDC fuer cosign keyless
      attestations: write
\end{lstlisting}

\begin{lstlisting}[language=bash,caption={Verifikation vor Deploy (\texttt{cd.yml})},label={lst:verify}]
cosign verify \
  --certificate-identity-regexp "https://github.com/$REPO/.github/workflows/.+" \
  --certificate-oidc-issuer "https://token.actions.githubusercontent.com" \
  "$IMAGE"

cosign verify-attestation --type slsaprovenance \
  --certificate-identity-regexp "https://github.com/$REPO/.github/workflows/.+" \
  --certificate-oidc-issuer "https://token.actions.githubusercontent.com" \
  "$IMAGE"
\end{lstlisting}

\subsection{Mapping Bedrohung $\to$ Maßnahme}
\begin{longtable}{p{5.0cm} p{9.5cm}}
\toprule
\textbf{Bedrohung (Tab.~\ref{tab:stride})} & \textbf{Maßnahme} \\
\midrule \endhead
Maintainer-Credential-Diebstahl & MFA, signierte Commits, CODEOWNERS-Review \\
Manipulation Pipeline-Definition & CODEOWNERS auf \texttt{.github/}, Required Checks \\
Bösartige Dependency & pip-audit, Dependabot, Pinning, SBOM-Diff \\
Base-Image-Replacement & Pinning + Trivy + Image-Digest in Deploy \\
Zu weite \texttt{permissions} & default-deny + per-Job opt-in \\
\texttt{pull\_request\_target} Leak & Trigger nicht verwendet; PRs ohne Secrets \\
Tag-Hijack in Registry & Cosign-Verify + Digest-Pinning beim Deploy \\
Fehlende Provenance & SLSA-Provenance via \texttt{attest-build-provenance} \\
Long-lived Cloud-Creds & OIDC-Federation, keine statischen Secrets \\
Fehlendes Approval Production & GitHub Environment Reviewer + Wait-Timer \\
\bottomrule
\caption{Bedrohungs--Maßnahmen-Matrix.}
\end{longtable}

\subsection{Fallstudie: \texttt{tj-actions/changed-files} (März 2025)}
\label{sec:tjactions}
Die Action \texttt{tj-actions/changed-files} wurde im März 2025
kompromittiert: Angreifer überschrieben mehrere veröffentlichte Tags so,
dass sie auf einen manipulierten Commit zeigten, der Secrets aus dem
Runner-Memory in die Build-Logs exfiltrierte. Workflows, die die Action
ohne SHA-Pinning eingebunden hatten (\texttt{tj-actions/changed-files@v44}),
zogen automatisch die manipulierte Version. Über zehntausend Repositories
waren betroffen, in mehreren Fällen wurden GitHub-PATs und Cloud-Credentials
öffentlich auf GitHub-Logs sichtbar.

Der Vorfall illustriert exakt die Bedrohungsklasse \emph{Source/Build-T}
aus Tabelle~\ref{tab:stride}. Die in dieser Arbeit umgesetzte Maßnahme --
durchgängiges Pinning aller Actions auf den Commit-SHA mit einem
nachgestellten Versions-Kommentar -- hätte den Angriff vollständig
neutralisiert: Selbst wenn der Tag \texttt{v44} auf einen bösartigen
Commit umgehängt wird, verweist der Pin
\texttt{tj-actions/changed-files@<sha>~\#~v44} weiterhin auf die
auditierte Version. Sigstore/Cosign hätte die Auswirkung zusätzlich
begrenzt, da das resultierende Artefakt keine gültige Signatur unter
der OIDC-Identität des Workflows getragen hätte.

\newpage

\section{Evaluation}\label{sec:evaluation}

\subsection{SLSA-Einstufung}\label{sec:slsa}
\emph{Supply-chain Levels for Software Artifacts} (SLSA) definiert in
Version~1.0~\cite{slsa} eine ordinale Skala für die Vertrauenswürdigkeit
der Build-Plattform und des Quellcode-Prozesses. Die Skala besteht aus
zwei orthogonalen Tracks (\emph{Build} und \emph{Source}); die jeweiligen
Levels sind kumulativ, d.\,h.\ jedes höhere Level fordert die
Anforderungen aller niedrigeren mit. Im Folgenden werden die Levels
zusammengefasst, mit der konkreten Umsetzung im untersuchten System
abgeglichen und die Begründung für \emph{nicht} erreichte höhere Stufen
transparent gemacht.

\paragraph{Build-Track.}
\begin{description}[leftmargin=1.4em,style=nextline]
  \item[L0] Keine besonderen Anforderungen. Der Build kann lokal,
        manuell, ohne jede Provenance erfolgen.
  \item[L1] Provenance existiert in irgendeiner Form -- der Build ist
        also reproduzierbar dokumentiert. Inhalt und Authentizität der
        Provenance werden noch nicht abgesichert.
  \item[L2] Provenance ist \emph{signiert} und wird von einem
        \emph{gehosteten Build-Service} erzeugt. Der Build läuft also
        nicht mehr auf einer Entwickler-Workstation; der
        Build-Service kann Provenance an seine eigene Identität
        binden.
  \item[L3] Der Build-Service garantiert zusätzlich
        \emph{Non-Forgeable Provenance} und \emph{Isolation} zwischen
        Builds: Provenance kann nicht durch den Build-Schritt selbst
        manipuliert werden, und ein Build kann nicht mit Artefakten
        oder Geheimnissen vorheriger Builds interagieren.
\end{description}

Die hier umgesetzte Pipeline erreicht \textbf{Build~L3}. Konkret:

\begin{itemize}[leftmargin=1.4em]
  \item GitHub-hosted Runner (ubuntu-24.04) ist ein gehosteter Service;
        jeder Job läuft in einer frisch provisionierten VM, was die
        Isolation\-sanforderung erfüllt (\emph{ephemeral builder}).
  \item Provenance wird mit \texttt{actions/attest-build-provenance}
        erzeugt und über die OIDC-Identität des Workflows signiert; das
        Signaturmaterial liegt nie im für die Build-Schritte
        zugreifbaren Speicher (Non-Forgeable).
  \item Das Build-Skript (\texttt{Dockerfile} und Workflow) ist
        versioniert, mit Branch-Protection und CODEOWNERS-Review
        geschützt.
\end{itemize}

\newpage

\paragraph{Source-Track.}
\begin{description}[leftmargin=1.4em,style=nextline]
  \item[L1] Quellcode ist versioniert (z.\,B.\ Git); Änderungen sind
        rückverfolgbar.
  \item[L2] Zusätzlich \emph{verifizierte Geschichte}: Commits sind
        kryptographisch signiert; der Repository-Provider verifiziert
        Identitäten beim Push.
  \item[L3] Zusätzlich \emph{Two-person review} und Schutz vor
        Force-Push / History-Rewrite.
\end{description}

Die Pipeline erreicht \textbf{Source~L3}: \texttt{main} ist
branch-protected mit CODEOWNERS-Review (Vier-Augen), Force-Push und
History-Rewrite sind serverseitig deaktiviert, signierte Commits sind
Pflicht-Status-Check.

\paragraph{Warum nicht höher?}
SLSA~v1.0 definiert offiziell nur Build-L0 bis L3 und Source-L1 bis
L3; es gibt also kein \enquote{L4}, das hier zu erreichen wäre.
Konzeptionell denkbare \emph{strenger als L3}-Eigenschaften wären
\emph{hermetische Builds} (kein Netzwerkzugriff zur Build-Zeit, alle
Inputs hash-gepinnt -- die Standard-pip-Installs verletzen das),
\emph{reproduzierbare Builds} (bit-identische Artefakte, bei
Container-Images nicht trivial wegen Zeitstempeln und Layer-IDs)
sowie \emph{Multi-Party-Builds} mit zwei unabhängigen Plattformen.
Diese Eigenschaften sind außerhalb der Reichweite eines Setups auf
einem einzigen Provider und werden im Ausblick (Kapitel~\ref{sec:fazit})
als weiterführende Forschungsrichtungen benannt. Die kombinierte
Erreichung von Build-L3 + Source-L3 ist mit handelsüblichen Mitteln
(GitHub Actions + Sigstore) ohne kommerzielle Add-ons darstellbar
und für den Studienarbeit-Scope ein bewusst gesetzter Höchststand.

\paragraph{Verifikation der Einstufung.}
Die Selbst-Einstufung wird durch externe Werkzeuge gegengeprüft:
OpenSSF Scorecard (Abschnitt~\ref{sec:scorecard}) bewertet einzelne
SLSA-relevante Checks objektiv;
\texttt{actions/attest-build-provenance} produziert maschinenlesbare
Provenance, die mit \texttt{slsa-verifier} gegen die deklarierten
Anforderungen geprüft werden kann.

\subsection{OpenSSF Scorecard}\label{sec:scorecard}
Der Scorecard-Workflow läuft wöchentlich. Erwartete Scores $\geq$~8/10 für
\emph{Branch-Protection}, \emph{Pinned-Dependencies}, \emph{Token-Permissions},
\emph{Signed-Releases}, \emph{SAST}, \emph{Vulnerabilities}.

\newpage

\subsection{Coverage der STRIDE-Kategorien}
Tabelle~\ref{tab:coverage} fasst zusammen, welche STRIDE-Kategorien durch die
Maßnahmen mitigiert sind.

\begin{table}[h]
\centering
\begin{tabular}{lcl}
\toprule
\textbf{Kategorie} & \textbf{Mitigiert} & \textbf{Wesentliche Maßnahme} \\
\midrule
Spoofing & vollständig & OIDC, signierte Commits, Cosign \\
Tampering & weitgehend & Pinning, Provenance, Verify-Gate \\
Repudiation & vollständig & Provenance + Rekor-Transparency-Log \\
Information Disclosure & weitgehend & gitleaks, least-privilege Tokens \\
Denial of Service & teilweise & concurrency-cancel, Cache-Strategie \\
Elevation of Privilege & vollständig & default-deny permissions, Environments \\
\bottomrule
\end{tabular}
\caption{STRIDE-Coverage durch implementierte Maßnahmen.}
\label{tab:coverage}
\end{table}

\subsection{Validierung gegen OWASP Top 10 CI/CD}\label{sec:owasp}
Zur Plausibilisierung der STRIDE-Analyse wird die abgeleitete
Maßnahmenliste zusätzlich gegen die OWASP-Top-10-CI/CD-Risiken
\cite{owaspcicd} geprüft. Tabelle~\ref{tab:owasp} zeigt die Abdeckung.

\begin{longtable}{p{0.9cm} p{4.8cm} p{8.2cm}}
\toprule
\textbf{Nr.} & \textbf{OWASP-Risiko} & \textbf{Wesentliche Maßnahme} \\
\midrule \endhead
CICD-1 & Insufficient Flow Control Mechanisms & Branch-Protection, CODEOWNERS, Environments mit Approval \\
CICD-2 & Inadequate Identity \& Access Mgmt    & OIDC-Federation, least-privilege \texttt{permissions} \\
CICD-3 & Dependency Chain Abuse                & Pinning, pip-audit, Dependabot, SBOM-Diff \\
CICD-4 & Poisoned Pipeline Execution (PPE)     & Kein \texttt{pull\_request\_target}, CODEOWNERS auf \texttt{.github/} \\
CICD-5 & Insufficient PBAC                     & Per-Job \texttt{permissions}, Environment-Reviewer \\
CICD-6 & Insufficient Credential Hygiene       & gitleaks, Push-Protection, Cosign keyless (keine Long-lived Keys) \\
CICD-7 & Insecure System Configuration         & Scorecard wöchentlich, Action-Pinning, default-deny Permissions \\
CICD-8 & Ungoverned Usage of 3rd-Party Services & Kuratierte Action-Liste, SHA-Pin, Renovate/Dependabot Audit \\
CICD-9 & Improper Artifact Integrity Validation & Cosign-Verify + SLSA-Provenance vor Deploy, Digest-Pinning \\
CICD-10 & Insufficient Logging \& Visibility   & SARIF-Upload, Rekor Transparency-Log, Audit-Logs (org-level) \\
\bottomrule
\caption{Abdeckung der OWASP-Top-10-CI/CD-Risiken durch die
implementierte Pipeline.}
\label{tab:owasp}
\end{longtable}

Alle zehn Risikoklassen werden adressiert. CICD-10 ist organisatorisch
abhängig (Audit-Logs auf Org-Ebene erfordern entsprechende
GitHub-Enterprise-Konfiguration) und damit nur teilweise innerhalb der
Pipeline selbst mitigierbar.

\subsection{KPIs und Sicherheitsmetriken}\label{sec:kpis}
Die bisherigen Evaluationsschritte (SLSA, Scorecard, STRIDE-Coverage,
OWASP-Mapping) bewerten den \emph{strukturellen} Reifegrad der
Pipeline. Im laufenden Betrieb muss zusätzlich beobachtet werden, ob
diese Mechanismen tatsächlich \emph{wirken} und welche Reibung sie
erzeugen. Die in Tabelle~\ref{tab:kpis} aufgeführten Kennzahlen
orientieren sich am SE-Framework-Quadranten \emph{Assurance}
(Tab.~\ref{tab:seframework}) und an den DORA-Metriken, ergänzt um
sicherheits\-spezifische Größen. Sie lassen sich in drei Cluster
gruppieren: \emph{Reaktionszeiten} (MTTD, MTTR pro Schweregrad)
messen die Wirksamkeit der Erkennungs-Pipeline; \emph{Coverage}
(Anteil signierter Artefakte, Branch-Protection-Verbreitung,
SAST/SCA-Abdeckung) misst die Verbreitung der Schutzmaßnahmen über
das Portfolio; \emph{Reibungsindikatoren} (False-Positive-Rate,
Pipeline-Bypass-Rate, Median-Build-Zeit) erfassen die in
Abschnitt~\ref{sec:humanfactor} diskutierte Anreizlage quantitativ
-- die False-Positive-Rate ist dabei der wichtigste Frühindikator,
ab ca.~20\,\% ist mit systematischem Umgehen der Gates zu rechnen.
Ergänzend bietet sich die \emph{Security Change Failure Rate} an --
der Anteil der Deployments, die wegen eines Sicherheits\-vorfalls
(nicht eines Funktions\-defekts) zurückgerollt werden müssen. Eine
vollständige Org-weite KPI-Erhebung erfordert GitHub Enterprise
Audit Log oder ein dediziertes SIEM und liegt außerhalb des in
Abschnitt~1.3 definierten Scopes.

\begin{table}[h]
\centering
\small
\begin{tabular}{p{4.6cm} p{4.4cm} p{4.4cm}}
\toprule
\textbf{Kennzahl} & \textbf{Datenquelle} & \textbf{Zielwert} \\
\midrule
MTTD (CRITICAL/HIGH) & pip-audit, Trivy, Dependabot & < 24\,h \\
MTTR (CRITICAL) & GitHub Security Alerts & < 7\,d \\
Anteil signierter Images & Cosign / Rekor & 100\,\% (Prod) \\
Branch-Protection Coverage & GitHub Org-API & 100\,\% (Prod-Repos) \\
False-Positive-Rate Gates & Build-Logs + Annotation & < 20\,\% \\
Pipeline-Bypass-Rate & Branch-API & 0\,\% \\
Security Change Failure Rate & Deploy-Logs + Incident-Tracker & $\to$ 0 \\
\bottomrule
\end{tabular}
\caption{Kennzahlen zur kontinuierlichen Bewertung der
Pipeline-Sicherheit.}
\label{tab:kpis}
\end{table}

\subsection{Compliance und regulatorischer Rahmen}\label{sec:compliance}
Die in dieser Arbeit umgesetzten Mechanismen erfüllen nicht nur
selbstgesetzte Sicherheitsziele, sondern adressieren zugleich externe
Anforderungen aus etablierten Standards und EU-Regulierung.
Tabelle~\ref{tab:compliance} ordnet die wichtigsten Rahmen den
relevanten Pipeline-Maßnahmen zu, ohne ein vollständiges Audit-Mapping
zu liefern.

\begin{table}[h]
\centering
\small
\begin{tabular}{p{3.0cm} p{4.4cm} p{6.0cm}}
\toprule
\textbf{Rahmen} & \textbf{Relevante Anforderung} & \textbf{Pipeline-Antwort} \\
\midrule
NIST SSDF\newline (SP\,800-218)~\cite{ssdf} &
PS (Protect Software), PW (Produce Well-Secured), RV (Respond to Vulnerabilities); insb.\ PW.4 (\enquote{Reuse Existing}) &
Signierte Artefakte + Provenance (PS); SAST/SCA + Threat-Model in Kap.~\ref{sec:threatmodel} (PW); Dependabot + SECURITY.md (RV); SHA-Pinning (PW.4) \\
\addlinespace
EU CRA~\cite{cra}\newline (ab 2027) &
SBOM-Pflicht; Schwachstellen-Handling über Produkt-Lebensdauer; 24h-Reporting bei aktiv ausgenutzten CVEs &
SBOM-Lebenszyklus (Abschn.~\ref{sec:sbomlifecycle}) erfüllt SBOM-Pflicht \emph{by construction}; Reporting flankiert durch SECURITY.md + Private Vulnerability Reporting \\
\addlinespace
NIS2 (EU 2022/2555) &
Art.\,21 Abs.\,2 lit.\,d (Lieferketten-Sicherheit); lit.\,e (Sicherheit bei Entwicklung und Wartung) &
Provenance, signierte Artefakte, Dependency-Tracking; voraussichtlich wirkungsstärkster regulatorischer Treiber für DevSecOps in DE \\
\addlinespace
ISO/IEC 27001:2022 &
Annex A.8.28 (Secure coding); A.8.29 (Security testing) &
Operationalisiert durch SAST/SCA und Coverage-Gate \\
\addlinespace
BSI IT-Grundschutz &
CON.8 Software-Entwicklung &
Kontrollierter Build-Prozess, Schwachstellen\-mgmt., Schutz der Entwicklungs\-umgebung sind durch die Pipeline umgesetzt \\
\bottomrule
\end{tabular}
\caption{Mapping der Pipeline-Maßnahmen auf Standards und Regulierung.}
\label{tab:compliance}
\end{table}

Eine vollständige Zertifizierung gegen diese Standards erfordert
allerdings organisatorische Belege jenseits des Pipeline-Scopes
(Schulungs\-nachweise, dokumentierte Verantwortlichkeiten,
Audit-Reports) und ist nicht Gegenstand dieser Arbeit. Die Pipeline
trifft die technischen Anforderungen ohne Sondermaßnahmen, weil
Provenance, Signatur und Vulnerability-Tracking zugleich Sicherheits-
und Compliance-Mechanismen sind -- eine Doppelfunktion, die im
Quadranten \emph{Assurance} (Tab.~\ref{tab:seframework}) verortet ist.

\subsection{Restrisiken}\label{sec:restrisks}
\begin{itemize}[leftmargin=1.4em]
  \item \textbf{Plattform-Kompromittierung} (GitHub, Sigstore): nicht durch
        die Pipeline mitigierbar; Vertrauensanker.
  \item \textbf{Insider mit Admin-Rechten} kann Branch-Protection deaktivieren.
        Mitigation: Audit-Logging, Org-Policies.
  \item \textbf{Zero-Day in transitiver Abhängigkeit} -- pip-audit erkennt nur
        bekannte CVEs.
  \item \textbf{Secret-Exfiltration via Build-Logs} bei fahrlässigem
        \texttt{echo}; nur durch Code-Review und \texttt{add-mask} mitigierbar.
\end{itemize}

\subsection{Zielkonflikt Geschwindigkeit vs.\ Assurance}
Strikte Gates (HIGH/CRITICAL fail, Coverage~80\,\%, manuelle Approvals)
verlangsamen Deployments. Die Pipeline mildert dies durch parallele Jobs,
GHA-Caches und \texttt{concurrency: cancel-in-progress} für PRs. Production
behält bewusst eine Approval-Hürde -- Geschwindigkeit darf hier nicht über
Assurance dominieren.

\subsection{Faktor Mensch und Anreize}\label{sec:humanfactor}
Technische Mechanismen entfalten ihre Wirkung nur, wenn die
beteiligten Akteure -- Entwickler, Reviewer, Plattform-Administratoren --
sie auch konsequent anwenden. Anderson formuliert im Security
Engineering Framework die These, dass Sicherheit nicht primär an
fehlenden Mechanismen scheitert, sondern an fehlerhaft gesetzten
\emph{Anreizen}: Die Person, die schützen \emph{soll}, hat oft nicht
denselben Nutzen vom erfolgreichen Schutz wie die Person, die im
Schadensfall den Verlust trägt. Dieser Anreizbruch ist im
DevSecOps-Kontext besonders prägnant, weil das Liefer- und das
Sicherheits-Ziel dieselbe Pipeline teilen. Die folgenden
Spannungsfelder sind in der Praxis besonders relevant.

\paragraph{Reibung und Anreize.}
Sicherheits-Gates wie Coverage-Schwellen oder Manual-Approval-Schritte
verlangsamen den Feedback-Zyklus. Wird die Reibung als
unverhältnismäßig empfunden, neigen Teams zum Umgehen --
\texttt{[skip ci]}-Commits, \texttt{--no-verify} bei pre-commit-Hooks,
pauschales Akzeptieren von \texttt{HIGH}-Findings als
\enquote{false positive}. Verhaltensökonomisch deckt sich das mit
dem in der Vorlesung diskutierten Modell der \emph{hyperbolischen
Diskontierung}: dem unmittelbaren Nutzen (PR gemerged) wird ein
unverhältnismäßig hoher Wert beigemessen gegenüber dem zukünftigen
Nutzen einer vermiedenen CVE-Ausnutzung. Die Gegenmaßnahmen wirken
strukturell, nicht moralisch: parallele Jobs +
\texttt{cancel-in-progress} kürzen Wartezeit; lokale
\texttt{pre-commit}-Hooks mit denselben Tools wie die CI fangen
Verstöße bereits vor dem Push -- der Schutz wird damit zur
\emph{schnellsten} Option, der Diskontierungseffekt arbeitet
zugunsten der Sicherheit. Die Trennung zwischen Build-Breakern
(\texttt{HIGH,CRITICAL}) und reinen SARIF-Reportings adressiert
zugleich \emph{alert fatigue}: Build-Abbrüche bleiben
handlungsrelevant, weniger kritische Findings verschwinden nicht,
erzwingen aber keinen Stillstand. Auf der Update-Seite wirkt
Dependabot analog: kleine, isolierte PRs durch die volle Pipeline
machen Aktualisieren zur \emph{billigsten} statt teuersten Option --
ein klassischer Anreiz-Hebel im Sinne des
SE-Frameworks (Tab.~\ref{tab:seframework}).

\paragraph{Verantwortungsverteilung und Shift Left.}
\enquote{Security ist Aufgabe des Security-Teams} ist ein verbreitetes
Anti-Pattern; \enquote{shift~left} verschiebt diese Verantwortung
explizit zu den Entwicklern. Intendierter Effekt mit Kehrseite:
Entwickler sollen über Findings entscheiden, für die sie weder
ausgebildet sind noch Zeit zugewiesen bekommen -- Standardreaktion
ist das Stummschalten per \texttt{\# nosec}. CODEOWNERS macht Reviews
für sicherheits\-kritische Pfade (\texttt{.github/},
\texttt{Dockerfile}) zum Default und delegiert das Reviewer-Recht an
die fachlich nächsten Personen; SARIF-Upload sorgt dafür, dass
Findings auch für andere Reviewer sichtbar bleiben. Die
verbleibende Schulungs- und Begleitarbeit ist eine organisatorische
Investition, die kein Tool ersetzt.

\paragraph{Falsche Sicherheit durch Zertifikate.}
Eine grüne Pipeline und ein hoher Scorecard-Wert können den Eindruck
erwecken, ein Repository sei \enquote{sicher} -- in der Vorlesung
behandelt unter den Stichworten \emph{Security Theatre} (Schneier)
und \emph{Datenschutzparadoxon} (Lücke zwischen geäußerter Präferenz
und tatsächlichem Verhalten). Auf Team-Ebene zeigt sich dasselbe
Muster: In Reviews wird Sicherheit als Wert benannt, im Tagesgeschäft
setzen sich Geschwindigkeit und Feature-Lieferung durch, sobald das
nominale Schutzniveau \emph{erscheint} ausreichend zu sein.

% ────────────────────────────────────────────────────────
\section{Fazit und Ausblick}\label{sec:fazit}

Die Arbeit zeigt, dass sich eine moderne CI/CD-Pipeline mit
zugänglichen Mitteln (GitHub Actions, Sigstore, Trivy, Syft) systematisch
entlang einer STRIDE-Analyse absichern lässt und dabei nachweisbare
Eigenschaften gemäß SLSA erreicht. Wesentlich ist nicht das einzelne Tool,
sondern die Kombination: \emph{least-privilege Identitäten, gepinnte
Abhängigkeiten, signierte Artefakte mit Provenance, verifizierte Deployments}.

Weiterführende Arbeiten könnten:
\begin{itemize}[leftmargin=1.4em]
  \item die Übertragbarkeit der Maßnahmen auf GitLab~CI quantitativ
        vergleichen,
  \item Policy-as-Code (z.\,B.\ Kyverno, Sigstore Policy Controller im Cluster)
        als zusätzliche Verifikationsschicht ergänzen,
  \item formale Analysen zu Pipeline-Konfigurationen (z.\,B.\ \texttt{zizmor},
        \texttt{actionlint}) mit Beweisbarkeit der \emph{permissions}-Minimalität
        verbinden.
\end{itemize}

\newpage

% ────────────────────────────────────────────────────────
\begin{thebibliography}{9}
\bibitem{slsa} OpenSSF, \emph{Supply-chain Levels for Software Artifacts (SLSA) v1.0}, \url{https://slsa.dev/spec/v1.0/}, abgerufen 2026.
\bibitem{ssdf} NIST, \emph{Secure Software Development Framework (SSDF) Version 1.1}, SP\,800-218, 2022.
\bibitem{stride} Microsoft, \emph{The STRIDE Threat Model}, Microsoft Learn, abgerufen 2026.
\bibitem{scorecard} OpenSSF, \emph{Scorecard}, \url{https://github.com/ossf/scorecard}, abgerufen 2026.
\bibitem{sigstore} Sigstore Project, \emph{Cosign / Fulcio / Rekor}, \url{https://www.sigstore.dev/}, abgerufen 2026.
\bibitem{owaspcicd} OWASP, \emph{Top 10 CI/CD Security Risks}, \url{https://owasp.org/www-project-top-10-ci-cd-security-risks/}, abgerufen 2026.
\bibitem{cra} Europäische Union, \emph{Cyber Resilience Act} (Verordnung (EU) 2024/2847), \url{https://eur-lex.europa.eu/eli/reg/2024/2847/oj}, 2024.
\end{thebibliography}

\end{document}